\section{Native tree-Prikry proofs}
\label{annex:p06-tree-prikry-forcing}

All tree operations and color choices are performed in the standing
ground \(M\). Greater conditions are stronger until an explicit host
subscript is supplied.
The host recursion, density, truth and ZFC inputs are exactly
\cref{thm:p06-ftr-recursion,thm:p06-ftr-density,thm:p06-ftr-truth,thm:p06-ftr-zfc}.

\begin{annexitem}{annex:p06.tree-prikry.structure}{thm:p06-tree-prikry-structure}{Pointwise completeness and level grafts}
\begin{proof}
The complement of every fewer-than-\(\kappa\)-sized set is in \(U\),
by intersecting the complements of its singletons. Thus measure-one
sets have size \(\kappa\), and every tail is in \(U\). If
\(\operatorname{cf}(\kappa)<\kappa\), the complements of a cofinal
family of bounded intervals would have empty intersection, contrary to
completeness. Hence \(\kappa\) is regular. If
\(e:\kappa\to\mathcal P(\lambda)\) were injective for
\(\lambda<\kappa\), choose the \(U\)-large membership or nonmembership
fiber for each of the \(\lambda\) bits, and intersect them. Injectivity
makes that measure-one intersection have at most one point, impossible.
Thus \(2^\lambda<\kappa\); normality was not used.

For \(r\) at or above the trunk, the cone \((T)_r\) retains precisely
the nodes end-extending \(r\). Their successor sets are unchanged
from \(T\), and rooted intermediate prefixes are retained. Hence
it is a \(U\)-tree with trunk \(r\), even though shorter ancestors
have been discarded.

Suppose \(T\subseteq S\) and the respective trunks are \(t,s\).
Since \(t\in S\), it is comparable with \(s\). If it were strictly
shorter than \(s\), its immediate successors in \(S\), hence in
\(T\), would have at most one point, because every \(S\)-node is
comparable with \(s\). This contradicts \(\operatorname{Suc}_T(t)\in U\).
Therefore \(t\) end-extends \(s\). A tree has only one such trunk,
by the same argument. Inclusion is reflexive and transitive, and mutual
inclusion gives the same tree and trunk, proving antisymmetry. Equal-trunk
inclusion gives the direct order and its inclusion in the full order.

For same-trunk \(T_\xi\), \(\xi<\mu<\kappa\), their intersection
retains the trunk and its rooted prefix structure. At a retained node
\(\eta\), its successor set is exactly
\(\bigcap_{\xi<\mu}\operatorname{Suc}_{T_\xi}(\eta)\), in \(U\).
This proves the intersection assertion and gives bounds for all direct
strengthening chains of length less than \(\kappa\). For the empty
chain use the tree of all increasing finite sequences. In particular,
same-trunk conditions are compatible. There are \(\kappa\) finite
trunks, so an antichain has at most \(\kappa\) members, proving
\(\kappa^+\)-cc.

For the grafting assertion, fix a relative level \(n\), and for each
\(\eta\in\operatorname{Lev}_{|t|+n}(T)\) choose
a direct subtree \(W_\eta\subseteq(T)_\eta\) with trunk \(\eta\).
Keep the part of \(T\) through that level, and above it take the union
of the \(W_\eta\). A node below the target level keeps all its old
successors. At and above it, the node lies in the unique cone determined
by its initial segment at the target level, and uses that cone's
measure-one successors. Hence the graft is a \(U\)-tree. No intersection
of the \(\kappa\) many distinct cone choices is required.
\end{proof}
\end{annexitem}

\begin{annexitem}{annex:p06.tree-prikry.host}{thm:p06-tree-prikry-host}{Same strengthening sets, different written orientation}
\begin{proof}
The host relation \(q\leq_h p\) is precisely the native strengthening
relation \(q\geq p\); it is the inverse of the source's weak-to-strong
\(\leq\). The corresponding direct relations have the same equality.
The weakest native condition is therefore the greatest host condition.
A common native strengthening is exactly a common host lower bound,
so compatibility and antichains are identical. The statements
\((\forall p)(\exists q\in D)q\geq p\) and
\((\forall p)(\exists q\in D)q\leq_h p\) are literally equivalent.
Native stronger-directedness and weaker-closure of a filter become host
lower-directedness and upward closure, with no set of conditions changed.
Thus the same filters meet the same ground dense sets.

The name constructor uses pairs of earlier names and conditions, not
their written order orientation. Its rank recursion is identical in
both presentations, and evaluation tests membership of conditions in
the same filter. Thus induction on name rank gives identical evaluations.
Forcing's atomic recursions quantify over strengthening conditions and
over dense subsets below a condition; replacing native \(q\geq p\)
by host \(q\leq_h p\) changes neither set. Induction on the name-rank
pair gives equal atomic forcing values. Induction on the formula then
gives equal negation and quantifier clauses, since their strengthening
domains are also identical. This proves the forcing bridge before any
semantic comparison. The current recursive density, truth and ZFC
theorems therefore apply via the host order, without pretending that
the native glyph \(\geq\) meant \(\leq\) all along.
\end{proof}
\end{annexitem}

\begin{annexitem}{annex:p06.tree-prikry.level-colors}{thm:p06-tree-prikry-level-colors}{Backward colors at finite heights}
\begin{proof}
Fix a relative target level \(n\). Color nodes on that level by \(f\).
Working backward, if colors \(c\) have been assigned to the children
of a node \(\eta\), partition \(\operatorname{Suc}_T(\eta)\) into
the fewer-than-\(\kappa\) color fibers. Exactly one fiber is in \(U\):
none would contradict completeness, and two would contradict their
disjointness. Assign its color to \(\eta\) and keep those children.
Do this for every node at the preceding level, then continue backward
to \(t\). This is a finite recursion on heights, with set-indexed
ground choices, not an intersection of all node fibers.

Now keep paths from \(t\) that follow the chosen fibers through height
\(n\), and above each retained height-\(n\) node retain its original
cone. Every retained node uses a \(U\)-large successor set. The result
is a direct subtree whose height-\(n\) nodes all have the root's
propagated color, and thus have the same original \(f\)-value.
For \(n=0\) the trunk's level is a singleton and needs no pruning.
Construct one such subtree \(S_n\) for each \(n<\omega\) and intersect
them. The same-trunk intersection lemma applies because
\(\omega<\kappa\). Each level remains a subset of its homogeneous
level, proving the simultaneous result.
\end{proof}
\end{annexitem}

\begin{annexitem}{annex:p06.tree-prikry.decision}{thm:p06-tree-prikry-decision}{Frozen levels, sibling decisions and a minimal trunk}
\begin{proof}
Fix \(p=\langle t,T\rangle\) and a sentence \(\sigma\). Write
\({}^{0}\sigma=\sigma\), \({}^{1}\sigma=\neg\sigma\), as in the source.
Assume for contradiction that no direct extension of \(p\) decides it.
Set \(T^0=T\). At stage \(n\geq1\), work at each node \(\eta\)
on relative level \(n-1\) of \(T^{n-1}\). Partition its successor set
into \(X_0,X_1,X_2\): a successor \(\alpha\) is in \(X_i\),
\(i<2\), if a direct \(U\)-subtree of
\((T^{n-1})_{\eta\mathbin{\frown}\langle\alpha\rangle}\)
forces \({}^{i}\sigma\); it is in \(X_2\) if neither decision has such
a witness. The first two sets are disjoint, since same-trunk witness
trees have a common intersection strengthening and opposite decisions
are inconsistent. Exactly one \(X_i\) belongs to \(U\).

Keep that successor fiber at \(\eta\). If \(i<2\), choose for each
retained child its witnessing subtree and graft it above that child;
if \(i=2\), keep its previous cone. Perform these independent grafts
at every parent on the level, retaining all earlier nodes. The resulting
\(T^n\) is a direct subtree. It agrees with \(T^{n-1}\) through
relative level \(n-1\), and the new level \(n\) will not be deleted
at any later stage. If a retained child permits a direct decision, its
whole new cone forces that decision; all retained siblings have the
same decision. This is an actual ground recursion: forcing for a fixed
sentence is definable, the node family is a set, and Choice selects the
local witness trees.

Let \(S=\bigcap_{n<\omega}T^n\). Pointwise countable completeness
makes it a \(U\)-tree with trunk \(t\), and all frozen levels remain.
By the contradiction assumption \(\langle t,S\rangle\) does not
decide \(\sigma\). The recursive negation clause and decision density
give a full extension \(\langle r,R\rangle\) forcing \(\sigma\).
Choose such an extension with the number \(n\) of new trunk points
minimal. It cannot have \(n=0\), which would be a direct decision of
the original condition. Let \(r^-\) be its immediate predecessor.

At stage \(n\), the child \(r\) of \(r^-\) has a direct
\(\sigma\)-witness: its own rooted cone \((R)_r\) is contained in
the stage cone and is a direct strengthening of the present \(R\).
Since \(r\) survives, the selected successor fiber at \(r^-\)
must have been \(X_0\). Therefore each surviving sibling \(r'\)
has \(\langle r',(S)_{r'}\rangle\forces\sigma\), by the grafted
witness and monotonicity. Any extension of
\(\langle r^-,(S)_{r^-}\rangle\) with a longer trunk passes through
one of these siblings and forces \(\sigma\). An extension with the
same trunk can append a successor from its measure-one successor set
and also reach such a condition. Thus conditions forcing \(\sigma\)
are dense above this parent in the native sense. The source-host
bridge and recursive density closure give
\(\langle r^-,(S)_{r^-}\rangle\forces\sigma\), contradicting the
minimal \(n\). Hence a direct decision exists.

No normality, global diagonal intersection or classical Prikry lemma
appears in this construction. The countable intersection is taken at
each node, and the lowering of trunk length uses full-order density.
\end{proof}
\end{annexitem}

\begin{annexitem}{annex:p06.tree-prikry.strong-dense}{thm:p06-tree-prikry-strong-dense}{Witnessed cones followed by level homogeneity}
\begin{proof}
On the nodes of \(T\) above \(t\), set \(c(\eta)=1\) if some
direct subtree of \((T)_\eta\), with trunk \(\eta\), is a member
of \(D\), and set \(c(\eta)=0\) otherwise. Assign zero also to any
shorter ancestor below the trunk, making \(c:T\to2\) a total map.
Build a sequence of
direct trees by processing levels \(n=0,1,\ldots\). At a retained
node \(\eta\) on relative level \(n\) with \(c(\eta)=1\), choose
a direct \(D\)-witness inside the current cone and graft it there.
Such a witness still exists: intersect an original \(D\)-witness
with the current same-trunk cone; the intersection is stronger than
the witness and hence in the native open \(D\). At nodes with color
zero keep the current cone. All levels through \(n\) stay fixed at
that stage and forever thereafter.

Intersect the resulting countable sequence, obtaining \(S_0\).
Whenever \(\eta\in S_0\) has \(c(\eta)=1\), its entire cone
\(\langle\eta,(S_0)_\eta\rangle\) strengthens the selected witness
and belongs to \(D\). Apply the proved finite-level color lemma to
\(c\) on \(S_0\), obtaining a direct tree \(S\) with constant
color at each relative level. This again uses no normality.

Density gives a member \(\langle r,R\rangle\in D\) stronger than
\(\langle t,S\rangle\). Its trunk lies at some relative level
\(n\); its own rooted cone strengthens it, stays in the open \(D\),
and witnesses \(c(r)=1\). Hence every node
at that level has color one and its pure cone in \(S\) belongs to
\(D\). Every cone at a later level strengthens the cone at its
initial level-\(n\) node; openness puts it in \(D\) as well.
This proves the uniform-level claim independently of the decision lemma.
\end{proof}
\end{annexitem}

\begin{annexitem}{annex:p06.tree-prikry.branch}{thm:p06-tree-prikry-branch}{Cofinal branch requirements and filter recovery}
\begin{proof}
Generic directedness makes all its trunks comparable, so their union
is increasing with domain contained in \(\omega\). Trunk length at
least \(n\) is a dense requirement: append successors along the tree.
Likewise a trunk point above any \(\alpha<\kappa\) is dense, since
the successor set at a trunk is measure one and its intersection with
the tail above \(\alpha\) is nonempty. Thus the union has domain
\(\omega\) and is cofinal in \(\kappa\). Its graph has the ground
name consisting of \((\check{\langle n,\alpha\rangle},p)\) whenever
\(n\) is a trunk coordinate of \(p\) and that coordinate is \(\alpha\).
Ground regularity of \(\kappa\) makes the sequence new over \(M\).

For \(A\in U\) and \(p=\langle t,T\rangle\), keep precisely
the extensions of \(t\) whose new points lie in \(A\), along with
the old trunk part. At a retained node its successors are
\(\operatorname{Suc}_T(\eta)\cap A\in U\); thus this is a direct
tree extension. The corresponding conditions are dense. A generic
meeting one has all subsequent branch points in \(A\).

For \(\langle t,T\rangle\in G\), every later initial segment of
\(C\) belongs to \(T\). Choose a generic trunk at least as long as
that segment and then a common strengthening with \(\langle t,T\rangle\).
Its trunk belongs to \(T\), and the rooted prefix property supplies
the intervening initial segments. Thus \(G\subseteq G(C)\).
Conversely a condition in \(G(C)\) is compatible with each member
of \(G\). Their trunks are both initial in \(C\); take the longer
trunk, which belongs to both trees, and intersect their cones there.
The same-trunk intersection is a common strengthening. Genericity
for the dense set of conditions stronger than, or incompatible with,
the candidate therefore gives a generic strengthening of it, and
weaker closure puts the candidate in \(G\). Hence \(G=G(C)\).

The branch is definable from \(G\), and this displayed ground-poset
Separation definition recovers \(G\) from \(C\). Name evaluation by
rank recursion then yields \(M[G]=M[C]\). All these are statements
about the ground and its supplied generic extension, not an assertion
of ambient measurability or ambient cardinality of \(\kappa\).
\end{proof}
\end{annexitem}

\begin{annexitem}{annex:p06.tree-prikry.preservation}{thm:p06-tree-prikry-preservation}{Direct tree fusion and the three cardinal regimes}
\begin{proof}
If \(p\forces\dot X\subseteq\check\gamma\), \(\gamma<\kappa\),
recursively for each \(\alpha<\gamma\) choose a direct stronger tree
deciding \(\check\alpha\in\dot X\), using the proved tree decision
theorem. At a limit stage intersect the preceding fewer-than-\(\kappa\)
same-trunk trees. At the end their intersection gives \(q\geq^*p\)
deciding every membership bit. The set
\(Y=\{\alpha<\gamma:q\forces\check\alpha\in\dot X\}\) belongs
to \(M\) by fixed-formula forcing definability and Separation.

To justify \(q\forces\dot X=\check Y\), a stronger condition forcing
inequality would, by forced boundedness, the existential clause and
Extensionality, have a further condition witnessing a symmetric-difference
point in \(\check\gamma\). The check-name membership clause provides
a further strengthening identifying it with some \(\check\alpha\),
\(\alpha<\gamma\). That contradicts the fixed bit decision. Decision
density then forces equality. For \(\gamma=0\), boundedness already
gives equality with the empty set. Such ground-set decisions are dense
above any starting condition forcing the boundedness; their union with
conditions incompatible with it is globally dense. Genericity and the
truth lemma therefore give
\(\mathcal P(\gamma)^{M[G]}=\mathcal P(\gamma)^M\).
This argument used the tree decision theorem and pointwise completeness,
not the classical normal-measure theorem.

For a ground cardinal \(\lambda<\kappa\), any extension surjection
\(f:\xi\to\lambda\), \(\xi<\lambda\), has a graph which can be
coded by a ground injection of \(\xi\times\lambda\) into an ordinal
below \(\kappa\). No bounded subsets makes its code, hence \(f\),
a ground set, contradicting ground cardinality. For \(\lambda=\kappa\),
choose a ground cardinal \(\mu\) with \(|\xi|^M<\mu<\kappa\),
using the proved strong-limit property. A surjection onto \(\kappa\)
would give one onto \(\mu\) by replacing all values at least \(\mu\)
with zero. That smaller function has a bounded graph code and leads
to the same contradiction. The unbounded original graph is not incorrectly
treated as a bounded code.

For a ground cardinal \(\lambda>\kappa\), suppose \(p\) forces
an onto function \(\dot f:\check\xi\to\check\lambda\),
\(\xi<\lambda\). For each coordinate choose a maximal antichain
above \(p\) deciding its value as a ground ordinal. Such decisions
are dense by check-name membership, and Choice supplies the antichain.
Each has size at most \(\kappa\), by the tree chain condition.
Its set \(S_\alpha\) of possible values has that bound. Maximality
and forcing density imply that \(p\) forces each value to lie in its
corresponding \(S_\alpha\). The range is therefore contained in the
ground union of these sets, of size at most
\(\max\{|\xi|^M,\kappa\}<\lambda\), contradicting surjectivity.
Ground noncardinal ordinals remain noncardinal by their existing bijections;
the general forcing theorem adds no ordinals. This proves all cardinal
claims, without normality.

Let \(\delta\) be a nonzero ground limit ordinal and
\(\mu=\operatorname{cf}^M(\delta)\), with ground increasing cofinal
\(h:\mu\to\delta\). Its persistence makes extension cofinality at
most \(\mu\). If \(\mu<\kappa\), a shorter extension cofinal map
\(f:\nu\to\delta\), \(\nu<\mu\), would give a cofinal map into
\(\mu\) by assigning the least \(\eta\) with \(f(\alpha)<h(\eta)\).
That map has a bounded graph code and would be in \(M\), contradicting
ground regularity of \(\mu\). If \(\mu>\kappa\), the possible-value
antichain argument for a shorter cofinal map gives a ground range bound
of size at most \(\max\{|\nu|^M,\kappa\}<\mu\); such a ground set
is bounded in \(\delta\), again a contradiction. If \(\mu=\kappa\),
compose \(h\) with the generic branch \(C\); it is cofinal of length
\(\omega\), and no finite set is cofinal in a limit ordinal. Zero and
successor cofinalities are unchanged by their ordinal definitions.
\end{proof}
\end{annexitem}

\begin{annexitem}{annex:p06.tree-prikry.family-criterion}{thm:p06-tree-prikry-family-criterion}{Capturing every ground stem-indexed tree family}
\begin{proof}
First note that \(G(C)\) is a filter for any increasing \(\omega\)-sequence
\(C\): the full tree condition belongs to it; two members have a
common stronger condition by intersecting their cones at the longer
initial trunk; and a weaker condition contains the stronger tree and
its rooted initial path, so remains in \(G(C)\).

Suppose \(G(C)\) is generic and a ground family \(s\mapsto T_s\)
is given. The conditions \(\langle s,S\rangle\) with
\(S\subseteq T_s\) are dense: from \(\langle s,T\rangle\), intersect
\(T\) and \(T_s\) with the same trunk. A generic member of this
dense set places every subsequent initial segment of \(C\) in
\(T_s\), proving the family property.

Conversely suppose the family property holds. Given a ground native
dense open \(D\), apply the strong dense theorem at every finite
stem \(s\), starting with the entire increasing tree above \(s\).
Choose a ground \(U\)-tree \(T_s\) and integer \(n_s\) such that
every cone in \(T_s\) at relative depth at least \(n_s\) belongs
to \(D\). All choices are a ground family. The criterion gives an
initial \(s=C\mathbin{\upharpoonright}N\) after which the branch
is in \(T_s\). Let \(r=C\mathbin{\upharpoonright}(N+n_s)\).
Then \(\langle r,(T_s)_r\rangle\in D\), and the remaining branch
is in that cone, so the condition belongs to \(G(C)\). Thus \(G(C)\)
meets every ground dense open set. For arbitrary dense sets, use their
stronger closure and the filter's weaker closure. It is generic.

Taking \(T_s\) to be the tree whose points after \(s\) lie in a
fixed \(A\in U\) recovers eventual containment in \(A\) as a necessary
condition. The proof nowhere replaces an arbitrary tree family by a
single global measure-one set when normality has not been assumed.
\end{proof}
\end{annexitem}

\begin{annexitem}{annex:p06.tree-prikry.normal-comparison}{thm:p06-tree-prikry-normal-comparison}{Where normal diagonalization enters the dense comparison}
\begin{proof}
Now and only now assume \(U\) normal. For the classical host condition
\((s,A)\), let \(T(s,A)\) consist of extensions of \(s\) with every
added point in \(A\). It is a \(U\)-tree since its successor sets
are the tails of \(A\). Set \(\iota(s,A)=\langle s,T(s,A)\rangle\).
The classical extension clauses imply
\(q\leq p\Rightarrow\iota(q)\geq\iota(p)\). Conversely this tree
inclusion puts the longer trunk's new points in the old tail, and each
possible next point of its tail must be in the old tail; hence it gives
exactly the classical extension clauses. Equal trunks similarly give
the direct-order correspondence. Increasing enumeration supplies the
already proved bridge to the native finite-set presentation
\cref{thm:p06-prikry-order-bridge}.

If the images are compatible, a common tree trunk end-extends both
original stems and lies in both image trees. Thus the stems are comparable
and the points of the longer after the shorter lie in its required tail.
Intersecting the two tails above that longer stem gives a classical
common strengthening. The converse follows by mapping a classical
common strengthening. Thus incompatibility is preserved and reflected.

For density, fix \(\langle s,T\rangle\). At each finite node \(\eta\)
extending \(s\) in \(T\), assign its measure-one successor set;
assign \(\kappa\) at irrelevant finite stems. The normal finite-stem
diagonal lemma of \cref{thm:p06-prikry-measure-tools} gives \(A\in U\)
such that every later point of \(A\) lies in the assigned successor
set of each earlier finite stem. Trim \(A\) above \(s\), handling
the empty stem by the lemma's explicit empty-stem intersection. Induction
on the length of a stem addition proves \(T(s,A)\subseteq T\).
This is the direct dense image condition. It is precisely the normality
step missing from the unrestricted non-normal standing hypotheses.

For a classical generic filter \(H\), take the native weaker closure
of \(\iota[H]\). It is a filter by the order and incompatibility
properties. A ground dense open tree set pulls back to a dense classical
set: first strengthen an image into the tree dense set and then use
density of the image to strengthen it into another image. Hence the
closure is tree-generic. Conversely, a tree generic filter has a cofinal
subset of image conditions, by meeting the image dense set below each
of its members. Its preimage is a classical generic filter: classical
dense requirements have dense images in the tree forcing, and the same
cofinal image argument gives directedness. These two constructions are
inverse. The image preserves trunks, so their generic sequences agree.
Both filters are recovered from that sequence by their respective
already proved formulas, giving the same generated extension \(M[C]\).

Finally let \(C\) be eventually in every \(A\in U\), and let
\(s\mapsto T_s\) be a ground tree family. For every finite \(\eta\),
intersect \(\operatorname{Suc}_{T_s}(\eta)\) over the finitely many
initial segments \(s\) of \(\eta\) for which \(\eta\in T_s\);
use \(\kappa\) if there are none. These are measure-one sets.
Normal finite-stem diagonalization gives \(A\in U\) whose later
points obey each such finite-node requirement, including the empty
node. Choose \(N\) after which \(C\) lies in \(A\), and put
\(s=C\mathbin{\upharpoonright}N\). Starting at \(s\in T_s\),
induction on subsequent branch length shows that its next point lies
in \(\operatorname{Suc}_{T_s}(\eta)\). The branch is therefore in
\(T_s\) after \(s\), establishing the tree-family criterion. Without
the additional normality, this final diagonal argument is not asserted.
\end{proof}
\end{annexitem}
